\section{A prime plus a multiple of a prime cube or fourth power}
Lee and O'Clarey asked whether every integer greater than 24 is the sum of a prime and a positive multiple of a prime square. Their Theorem~1.3 establishes this under Dirichlet GRH. Campbell also reports a direct unconditional proof in the revision of 24 November 2026. Increasing the exponent to three gives the following different cutoff. \href{https://arxiv.org/html/2605.02426v1}{Lee and O'Clarey}; \href{https://arxiv.org/abs/2609.14714v2}{Campbell}.

\medskip\noindent\textbf{Theorem.} Exactly 62 positive integers fail to have a representation
\[
                  n=p+m r^3,\qquad p,r\text{ prime},\quad m\ge1.
\]
The largest exception is $10\,500$. The complete list is in \texttt{checks/cube.out}.

\smallskip\noindent\textit{Proof.}
Let $q$ be the smallest prime that does not divide $n$, and put $P_q=\prod_{r<q,\ r\ {\rm prime}}r$. Then $P_q\mid n$. If a prime $p<n$ satisfies $p\equiv n\pmod{q^3}$, it supplies the required representation. Corollary~1.2 of Lamzouri, Li and Soundararajan supplies such a congruent prime with
\[
                       p\le [\varphi(q^3)\log(q^3)]^2
\]
under Dirichlet GRH. This is the precise input supplied by the Lean theorem, together with the published implication. \href{https://arxiv.org/abs/1309.3595}{Lamzouri, Li and Soundararajan}.

For $q\ge43$ the displayed upper bound is less than $q^8$ since $3\log q<q$. Moreover, $P_q>q^8$ for all such primes. The checker compares integers for $43\le q\le257$. Thereafter Bertrand's postulate puts the next prime $s$ below $2q$, so $P_s=qP_q>q\cdot q^8>(2q)^8>s^8$. Consequently $p<n$ throughout this range.

For $q\le41$ let $M_q$ be the maximum of the least prime in each reduced class modulo $q^3$. A sieve through eight million supplies a representative of every class for each of these moduli. As soon as all $q^2(q-1)$ classes appear, the last prime used equals $M_q$. Every $n>M_q$ with this least nondivisor $q$ is covered. The remaining candidates are the positive multiples of $P_q$ up through $M_q$ that are coprime to $q$. For each candidate the checker visits every prime cube below $n$ until it finds a positive multiple with prime complement or exhausts the possibilities.
\[
\begin{array}{c|rrrrrrrr}
q&2&3&5&7&11&13&17&19\\ \hline
M_q&17&109&1621&9371&47737&120401&349529&540961\\
\text{candidates}&9&36&216&268&207&48&11&1\\
\text{exceptions}&6&11&25&15&5&0&0&0
\end{array}
\]
For $q=23,29,31,37,41$ there are no remaining candidates. This accounts for every positive integer. The finite calculation uses only integer arithmetic.

The last exception can also be checked by hand. For $n=10\,500$, a cube with base $r=2,3,5,7$ forces $p=r$, and $r^3\nmid n-r$ in each case. The other possible bases are $11,13,17,19$. In order of increasing positive multiplier $m$, proper divisors of all positive complements $n-mr^3$ are, respectively,
\[
(53,2,3,2,5,2,7),\qquad(19,2,3,2),\qquad(37,2),\qquad(11).
\]
This proves sharpness.\hfill$\square$
\subsection*{Fourth powers}
The same argument gives an exact cutoff when $r^3$ is replaced by $r^4$.
\textbf{Exactly 811 positive integers fail to be a prime plus a positive multiple of a prime fourth power. The largest is $1\,561\,560$.}

Retain $q$ and $P_q$ from the preceding proof. For $q\ge53$, the least congruent prime modulo $q^4$ is at most $[4\varphi(q^4)\log q]^2<q^{10}$. Exact comparisons show $P_q>q^{10}$ for all primes from 53 through 1031. Bertrand's postulate propagates the inequality from any $q>2^{10}$ to its successor. Thus $n$ exceeds the least congruent prime for the entire tail.

For $q\le23$, enumerate all reduced classes modulo $q^4$ as above, then test the remaining multiples of $P_q$. The corresponding maxima of the least primes, in increasing order of $q$, are
\[
 41,\ 647,\ 20063,\ 92369,\ 1189213,\ 2987429,\ 9951083,\ 19423919,\ 42026629.
\]
The nine groups yield, respectively, $12,48,230,255,212,46,7,1,0$ exceptions. When testing another prime base $r$ that divides $n$, the only possible prime complement is $r$ itself. This observation shortens the finite check.

For the six remaining primes $29\le q\le47$, we have $4\log q<16$. Hence the least-prime theorem covers $n>B_q=[16q^3(q-1)]^2$. For each positive multiple $n$ of $P_q$ through $B_q$ with $q\nmid n$, the checker exhibits a prime in the residue class of $n$ modulo $q^4$ strictly below $n$. The numbers of inputs in these six groups are $516676,30586,4130,256,9,0$. Sieve lookup and trial division certify all the primes. The full exceptional list and reproducible computation are in \texttt{checks/fourth.out} and \texttt{checks/fourth.py}.

As an independent finite check, for every prime $r$ with $r^4\le2\,000\,000$, one can take the first prime $p$ in every occupied class modulo $r^4$ and mark $p+r^4,p+2r^4,\ldots$ through two million. We obtained the identical list of 811 unmarked positive integers by this procedure. The primorial argument above certifies completeness beyond this finite check.
